\section{Introduction}

Security compliance is no longer confined to a small set of safety- or security-critical domains~\cite{moyon2020security-compliance-agile}.
Regulations such as the EU Cyber Resilience Act~\cite{eu2024cra} and a growing body of sector-specific standards (e.g., IEC 62443 for industrial automation and control system development~\cite{iec2018secure}) are making security a baseline expectation for nearly every organization that builds or operates software.
In organizations leveraging agile or DevOps, this circumstance requires engineering teams to continuously address security throughout the development process~\cite{fitzgerald2017continuous, turpe2017scrum}.
In these settings the product backlog is the central planning artifact~\cite{iso2018sse, ferreira2022writing-agile-requirements}, yet backlog items rarely carry explicit security or compliance annotations~\cite{turpe2017scrum}.
Security considerations are therefore mostly embedded implicitly in short, free-form engineering tasks. This leaves product teams without clear guidance on specific security matters and creates a practical mismatch between continuous software engineering and compliance efforts.
Security planning and reviewing is still often performed manually and outside the main backlog flow, which does not scale to fast development cycles~\cite{moyon2020security-compliance-agile}.
Moyón et al.~\cite{moyon2024challenges} describe this tension in secure continuous software engineering through several open practical challenges, such as making security architecture visible in backlog artifacts and documentation, involving security activities with minimal lead-time burden, and enabling security knowledge and ownership within engineering teams.


This manuscript reports on building, deploying, and evaluating a natural-language backlog enrichment system that addresses this gap directly in the artifact developers already use.
The system surfaces security-relevant context at the product planning stage by attaching security requirements relevant to the organization to flagged backlog items.
The work was developed and evaluated at a large enterprise in highly regulated domains. To enable reproduction, we additionally evaluate the approach on public security guidelines.

The paper makes three contributions:
\begin{itemize}
    \item[\textbf{C1}] An expert-annotated dataset of $288$ backlog items carrying binary security-relevance labels, drawn from two public Atlassian projects and multi-rated by $9$ security practitioners.
    \item[\textbf{C2}] A recall-oriented binary classifier that flags security-relevant backlog items as a first-stage filter for the backlog enrichment system.
    \item[\textbf{C3}] A retrieval pipeline that links backlog items to applicable security requirements from requirement documents.
\end{itemize}

We evaluate the three contributions from complementary perspectives: inter-rater reliability and dataset characterization for C1, in-distribution performance of C2 on C1 and cross-distribution transfer to the Wu et al.~\cite{wu2022data-quality} datasets, and a preliminary evaluation via two practitioner interviews on industrial backlogs for C3, which covers both proprietary and public security policy corpora.

Across the three contributions, a single overarching finding emerges. Security-relevance of backlog items can be evaluated based on the ordinary item text, and concrete security requirements can be attached to those items in a form that engineers and security experts can act on without leaving their existing workflows.
